\Article{4}
Die Aufgabe der Königlichen Gesellschaft der Wissenschaften verlangt, dass die Abbildung dem Abgebildeten in den kleinsten Teilen ähnlich sein soll. Zunächst ist diese Bedingung analytisch auszudrücken.

Durch Differentiation der Funktionen von \(t,u\), durch die \(x,\allowbreak y,\allowbreak z,\allowbreak X,\allowbreak Y,\allowbreak Z\) ausgedrückt werden, mögen die folgenden Gleichungen entstehen:
\[
\begin{aligned}
\mathrm{d}x&=a\,\mathrm{d}t+a'\,\mathrm{d}u\\
\mathrm{d}y&=b\,\mathrm{d}t+b'\,\mathrm{d}u\\
\mathrm{d}z&=c\,\mathrm{d}t+c'\,\mathrm{d}u\\
\mathrm{d}X&=A\,\mathrm{d}t+A'\,\mathrm{d}u\\
\mathrm{d}Y&=B\,\mathrm{d}t+B'\,\mathrm{d}u\\
\mathrm{d}Z&=C\,\mathrm{d}t+C'\,\mathrm{d}u
\end{aligned}
\]
\noindent Die vorgeschriebene Bedingung verlangt erstens, dass alle unendlich kleinen Linien, die von einem Punkt der ersten Fläche ausgehen und in ihr liegen, den entsprechenden Linien der zweiten Fläche proportional sind, und zweitens, dass die ersteren untereinander dieselben Winkel bilden wie die letzteren.

Ein solches Linienelement auf der ersten Fläche ist
\[
\begin{aligned}
={}&\BareRoot{\bigl((aa+bb+cc)\mathrm{d}t^2+2(aa'+bb'+cc')\mathrm{d}t\mathord{.}\mathrm{d}u}\\
&\qquad+(a'a'+b'b'+c'c')\mathrm{d}u^2\bigr)
\end{aligned}
\]
\noindent und das entsprechende auf der zweiten Fläche
\[
\begin{aligned}
={}&\BareRoot{\bigl((AA+BB+CC)\mathrm{d}t^2+2(AA'+BB'+CC')\mathrm{d}t\mathord{.}\mathrm{d}u}\\
&\qquad+(A'A'+B'B'+C'C')\mathrm{d}u^2\bigr).
\end{aligned}
\]
\noindent Sollen beide unabhängig von \(\mathrm{d}t\) und \(\mathrm{d}u\) in einem bestimmten Verhältnis zueinander stehen, so müssen offenbar die drei Größen
\[
aa+bb+cc,\quad aa'+bb'+cc',\quad a'a'+b'b'+c'c'
\]
\noindent jeweils den folgenden drei Größen proportional sein:
\SourceLeaf{200}
\[
AA+BB+CC,\quad AA'+BB'+CC',\quad A'A'+B'B'+C'C'
\]
\noindent Wenn den Endpunkten eines zweiten Elements auf der ersten Fläche die Werte
\[
t,u\quad\text{und}\quad t+\delta t,u+\delta u
\]
\noindent entsprechen, so ist der Kosinus des Winkels, den es mit dem ersten Element bildet,
\[
=\frac{\begin{gathered}(a\,\mathrm{d}t+a'\,\mathrm{d}u)(a\delta t+a'\delta u)+(b\,\mathrm{d}t+b'\,\mathrm{d}u)(b\delta t+b'\delta u)\\+(c\,\mathrm{d}t+c'\,\mathrm{d}u)(c\delta t+c'\delta u)\end{gathered}}{\begin{gathered}\BareRoot{((a\,\mathrm{d}t+a'\,\mathrm{d}u)^2+(b\,\mathrm{d}t+b'\,\mathrm{d}u)^2+(c\,\mathrm{d}t+c'\,\mathrm{d}u)^2)}\\\mathord{.}((a\delta t+a'\delta u)^2+(b\delta t+b'\delta u)^2+(c\delta t+c'\delta u)^2)\end{gathered}}
\]
\noindent und für den Kosinus des Winkels zwischen den entsprechenden Elementen auf der zweiten Fläche ergibt sich ein ganz ähnlicher Ausdruck, wenn nur \(a,b,c,a',b',c'\) durch \(A,B,C,A',B',C'\) ersetzt werden. Offenbar werden beide Ausdrücke gleich, wenn die obige Proportionalität gilt. Damit ist die zweite Bedingung bereits in der ersten enthalten, was auch bei einigem Nachdenken von selbst klar wird.

Der analytische Ausdruck der Bedingung unserer Aufgabe lautet demnach, dass
\[
\frac{AA+BB+CC}{aa+bb+cc}=\frac{AA'+BB'+CC'}{aa'+bb'+cc'}=\frac{A'A'+B'B'+C'C'}{a'a'+b'b'+c'c'}
\]
\noindent gelten muss. Dies wird eine endliche Funktion von \(t\) und \(u\) sein, die wir \(=mm\) setzen wollen. Dann drückt \(m\) das Verhältnis aus, in dem die linearen Größen auf der ersten Fläche in ihrer Abbildung auf der zweiten vergrößert oder verkleinert werden (je nachdem, ob \(m\) größer oder kleiner als 1 ist). Dieses Verhältnis wird im Allgemeinen von Stelle zu Stelle verschieden sein: Im besonderen Fall, dass \(m\) konstant ist, besteht auch in den endlichen Teilen eine vollkommene Ähnlichkeit, und wenn darüber hinaus \(m=1\) ist, besteht eine vollkommene Gleichheit, und die eine Fläche lässt sich auf die andere abwickeln.

